\section{Accept/Reject Methods}
\label{sec:2.2}

In this section, we study rejection sampling as a prototypical example of an \textit{accept/reject} method.
The idea is to sample from a \textit{proposal} (also known as reference or instrumental) distribution
and discard samples that are unlikely under the \textit{target} distribution of interest. We denote the
target distribution by $f$ and the proposal distribution by $g$. Typically, the target is expensive or
impossible to sample from directly, but the proposal is easy to sample from. We present rejection
sampling in Subsection~\ref{sec:2.2.1} and a related accept/reject approach for Approximate Bayesian
Computation (ABC) in Subsection~\ref{sec:2.2.2}.

\subsection{Rejection Sampling}
\label{sec:2.2.1}

Before introducing the rejection sampling algorithm, we start with a simple identity and an example, both of which will correspond to a uniform proposal distribution $g$ in the subsequent
developments.

First, the identity
\begin{equation}
  f(x) = \int_0^{f(x)} 1\, du
  \label{eq:2.5}
\end{equation}
shows that $f(x)$ is the marginal of a random vector that is uniformly distributed on the area under
the density $f(x)$, that is, on the region $\{(x,u) : 0 \le u \le f(x)\}$. In the next example, we use the
identity \eqref{eq:2.5} to obtain an algorithm to sample from a beta distribution.

\begin{figure}[htbp]
\centering
\includegraphics[width=0.75\textwidth]{figures/fig_02_3}
\caption{\textit{Rejection sampling for a Beta(2,2) distribution, see Example~\ref{ex:2.5}. The target p.d.f.\
is given by $f(x) = 6x(1-x)$ for $x \in (0,1)$. Proposed samples are obtained by uniformly sampling in
the rectangle with base 1 and height $M$. We accept those samples that lie below the curve of $f$. The first
coordinate of each accepted sample is distributed according to the target.}}
\label{fig:2.3}
\end{figure}

\begin{example}[Rejection Sampling: Beta Distribution]
\label{ex:2.5}
Let $f = \Beta(\alpha, \beta)$ so that $f(x) \propto
x^{\alpha-1}(1-x)^{\beta-1}$ for $x \in (0,1)$, and suppose that $\alpha, \beta > 1$. The identity \eqref{eq:2.5} suggests that to
obtain a draw $X^{(n)} \sim f$, we could draw uniformly in the rectangle with base 1 and height $M$
until we get a draw under the curve $f$ (see Figure~\ref{fig:2.3}), and then keep the first coordinate. The
following pseudocode formalizes the procedure:

\begin{enumerate}
  \item For $n = 1, \ldots, N$ do:
  \begin{enumerate}
    \item \textbf{Step 1:} Generate independently $Z^{(n)} \sim \Unif(0,1)$ and $U^{(n)} \sim \Unif(0,1)$, so $(Z^{(n)}, MU^{(n)})$ is uniform on the rectangle $[0,1] \times [0,M]$.
    \item \textbf{Step 2:} Set $X^{(n)} = Z^{(n)}$ if $f(Z^{(n)}) \ge MU^{(n)}$ (that is, if $(Z^{(n)}, MU^{(n)})$ lies under the curve $f$). Otherwise, go back to step 1.
  \end{enumerate}
\end{enumerate}

Any $X^{(n)}$ defined this way has p.d.f.\ $f$ as desired. Note that the conditional probability that
$(Z^{(n)}, MU^{(n)})$ is accepted if $Z^{(n)} = z$ is
\[
  \Prob\!\left(MU^{(n)} < f(Z^{(n)}) \,\Big|\, Z^{(n)} = z\right) = \Prob\!\left(U^{(n)} < \frac{f(z)}{M}\right) = \frac{f(z)}{M}.
\]
Since $U^{(n)}$ and $Z^{(n)}$ are independent, we can rewrite the above procedure as follows:

\medskip
\begin{enumerate}
  \item For $n = 1, \ldots, N$ do:
  \begin{enumerate}
    \item \textbf{Step 1:} Draw $Z^{(n)} \sim \Unif(0,1)$.
    \item \textbf{Step 2:} Set $X^{(n)} := Z^{(n)}$ with probability $\dfrac{f(Z^{(n)})}{M}$. Otherwise, go back to step~1.
  \end{enumerate}
\end{enumerate}
\medskip

\end{example}

Example~\ref{ex:2.5} uses that the density $\Beta(\alpha,\beta)$, $\alpha, \beta > 1$ is bounded within a (finite) rectangle.
The proposal distribution can then be taken to be uniform. We now generalize the procedure
to cases where the range or the domain of $f$ are unbounded. To that end, we use a proposal $g$
satisfying the following conditions:
\begin{itemize}
  \item $g$ can be sampled from and evaluated.
  \item $g$ has compatible support with $f$ (i.e.\ $g(x) > 0$ when $f(x) > 0$).
  \item There is $M > 0$ such that, for all $x \in E$, $f(x) \le Mg(x)$.
\end{itemize}

Notice that for the target $f = \Beta(2,2)$ in the example in Figure~\ref{fig:2.3}, the choices $g = \Unif(0,1)$ and
$M = 3/2$ satisfy these conditions. We are ready to introduce the rejection sampling algorithm:

\begin{algorithm}[H]
\caption{Rejection Sampling}
\label{alg:2.3}
\begin{algorithmic}[1]
\Require Target $f$, proposal $g$, constant $M > 0$ such that, for all $x \in E$, $f(x) \le Mg(x)$, sample size $N$.
\State \textbf{For} $n = 1, \ldots, N$ \textbf{do:}
\State \quad \textbf{Step 1:} Draw $Z^{(n)} \sim g$ independently of all other draws.
\State \quad \textbf{Step 2:} Set $X^{(n)} := Z^{(n)}$ with probability $\dfrac{f(Z^{(n)})}{Mg(Z^{(n)})}$. Otherwise, go to step~1.
\Ensure Sample $\{X^{(n)}\}_{n=1}^N \iid f$.
\end{algorithmic}
\end{algorithm}

For a convenient implementation of steps 1 and 2, we may proceed as follows:

\medskip
\begin{enumerate}
  \item \textbf{Step 1:} Draw $Z^{(n)} \sim g$ and $U^{(n)} \sim \Unif(0,1)$ independently.
  \item \textbf{Step 2:} Accept $X^{(n)} = Z^{(n)}$ if $U^{(n)} \le \dfrac{f(Z^{(n)})}{Mg(Z^{(n)})}$. Otherwise, go to step~1.
\end{enumerate}
\medskip

\noindent\fbox{\parbox{\dimexpr\linewidth-2\fboxsep-2\fboxrule}{%
\textbf{Pros and Cons}: Rejection sampling can be used when $f$ and $g$ are only known up to a normalizing constant (see Exercise~2.5). The method is very versatile and can be applied in general state space $E$, beyond the Euclidean setting considered here. Potential disadvantages include that (i) samples that are not accepted are thrown away; (ii) the time to obtain each sample is random; and (iii) the method suffers from the curse of dimension, as will be discussed below.}}
\medskip

The next result shows that the rejection sampling algorithm produces draws distributed according to the target. It also characterizes the acceptance rate under the assumption that $f$ and $g$ are normalized; otherwise, the second and third assertions need to be adjusted (see Exercise~2.5).

\begin{theorem}[Rejection Sampling Acceptance Rate]
\label{thm:2.6}
\begin{enumerate}
  \item The draws $X^{(n)}$ generated by the rejection sampling algorithm have distribution $f$ and are independent.
  \item The probability that a draw $Z^{(n)}$ generated in step 1 is accepted in step 2 is $1/M$.
  \item The algorithm will finish in finite time (with probability one). The expected number of iterations to accept each sample $X^{(n)}$ is $M$.
\end{enumerate}
\end{theorem}

\begin{proof}
First, the independence of the $X^{(n)}$ follows from the fact that the $Z^{(n)}$ are drawn independently. Let now $A \subset E$. Note that
\begin{align*}
  \Prob\bigl(Z^{(n)} \in A\ \&\ Z^{(n)} \text{ is accepted}\bigr)
  &= \int_A \frac{f(x)}{Mg(x)} g(x)\, dx \\
  &= \frac{\int_A f(x)\, dx}{M},
\end{align*}
where $\dfrac{f(x)}{Mg(x)}$ is the probability that $Z^{(n)}$ is accepted given that $Z^{(n)} = x$ and $g$ is the density of $Z^{(n)}$. Setting $A = E$ and noting that $f$ integrates to 1 gives
\[
  \Prob\bigl(Z^{(n)} \text{ is accepted}\bigr) = \frac{1}{M},
\]
proving the second claim in the statement. Now,
\begin{align*}
  \Prob(X^{(n)} \in A)
  &= \Prob\bigl(Z^{(n)} \in A \,|\, Z^{(n)} \text{ is accepted}\bigr) \\
  &= \frac{\Prob(Z^{(n)} \in A\ \&\ Z^{(n)} \text{ is accepted})}{\Prob(Z^{(n)} \text{ is accepted})} \\
  &= \frac{\dfrac{\int_A f(x)\,dx}{M}}{\dfrac{1}{M}} = \int_A f(x)\, dx,
\end{align*}
which proves the first claim. The third claim follows by noting that $M$ is the expected value of a geometric random variable with parameter $1/M$.
\end{proof}

\paragraph{Choice of Upper-Bound}
Theorem~\ref{thm:2.6} shows that if $f$ and $g$ are p.d.f.s such that $f(x) \le
Mg(x)$ for every $x \in E$, then Algorithm~\ref{alg:2.3} has acceptance rate $1/M$. Therefore, given $f$ and
$g$, $M$ should be chosen as small as possible while still satisfying the requirement that $f(x) \le
Mg(x)$ for every $x \in E$. This way the acceptance rate is maximized. Note that it always holds
that $M \ge 1$, since $1 = \int_E f(x)\, dx \le M \int_E g(x)\, dx = M$.

\paragraph{Choice of Proposal}
We want our proposal to be (1) easy to sample from; and (2) close to the
target, so that we can find a small constant $M \ge 1$ satisfying $f(x) \le Mg(x)$ for every $x \in E$.
These two goals are often in conflict with each other and a compromise must be found. In the
extreme case where $f = g$ one could choose $M = 1$ (lowest possible bound), which is clearly
not optimal if $f$ is extremely expensive (or impossible) to sample from directly! A common
strategy is to look for proposals in some parametric family $\{g_\theta, \theta \in \Theta\}$. If for each $\theta \in \Theta$ we
can find $M_\theta$ such that $f(x) \le M_\theta g_\theta(x)$ for every $x \in E$, then it is recommended to use as
proposal the density with parameter $\theta^*$ that minimizes $M_\theta$ over $\theta \in \Theta$.

\begin{remark}
\label{rem:2.7}
Note that if $E \subset \R^d$ is unbounded and $f(x) \le Mg(x)$ for every $x \in E$, then $g$
necessarily has fatter tails than $f$. Several Monte Carlo methods (e.g.\ importance sampling in
Chapter~3) share this principle that the proposal should have fatter tails than the target.
\end{remark}

\paragraph{Envelope Rejection Sampling}
If evaluating the target density $f$ is expensive, delaying the
acceptance can improve the efficiency. Specifically, let $\ell$ be a function that is cheap to evaluate
and such that $\ell(x) \le f(x)$ for all $x \in E$. Then, we can modify the rejection sampling algorithm
by delaying acceptance as follows:

\medskip
\begin{enumerate}
  \item \textbf{Step 1:} Sample $Z^{(n)} \sim g$, $U^{(n)} \sim \Unif(0,1)$.
  \item \textbf{Step 2:} Accept $X^{(n)} = Z^{(n)}$ if $U^{(n)} \le \dfrac{\ell(Z^{(n)})}{Mg(Z^{(n)})}$. Otherwise, go to step~3.
  \item \textbf{Step 3:} Accept $X^{(n)} = Z^{(n)}$ if $U^{(n)} \le \dfrac{f(Z^{(n)})}{Mg(Z^{(n)})}$. Otherwise, go to step~1.
\end{enumerate}
\medskip

\paragraph{Rejection Sampling and the Curse of Dimension}
Suppose that, for a given target $f$ and
proposal $g$, we have found a constant $M$ with $f(x) \le Mg(x)$ for every $x \in E = \R$. We know
that the probability of accepting a sample from $g$ in the rejection sampling algorithm is $\frac{1}{M}$. Now
consider a $d$-dimensional i.i.d.\ setting, where $f_d$ and $g_d$ are p.d.f.s on $E = \R^d$ defined by
\begin{align*}
  f_d(x) &:= f(x_1) \times \cdots \times f(x_d), \quad x = (x_1, \ldots, x_d) \in \R^d, \\
  g_d(x) &:= g(x_1) \times \cdots \times g(x_d), \quad x = (x_1, \ldots, x_d) \in \R^d.
\end{align*}
To implement rejection sampling with target $f_d$ and proposal $g_d$ we can only guarantee that, for
every $x \in \R^d$, the following upper bound holds:
\[
  \frac{f_d(x)}{g_d(x)} = \frac{f(x_1) \times \cdots \times f(x_d)}{g(x_1) \times \cdots \times g(x_d)} \le M^d.
\]
Thus, Theorem~\ref{thm:2.6} shows that the probability of accepting each sample from $g_d$ decreases exponentially with $d$. In other words, the expected number of iterations needed to obtain one sample
grows exponentially with $d$.

\subsection{Approximate Bayesian Computation}
\label{sec:2.2.2}

In Bayesian statistics, inference on a parameter $\theta$ given data $y$ is based on the posterior distribution $f(\theta|y)$. The posterior is obtained by combining a likelihood model $f(y|\theta)$ and a prior
density $f(\theta)$ using Bayes' formula:
\[
  f(\theta|y) = \frac{1}{c}\, f(y|\theta) f(\theta),
\]
where $c$ is a normalizing constant so that $f(\theta|y)$ integrates to 1. Approximate Bayesian Computation (ABC) refers to a family of methods designed to obtain \textit{approximate} posterior samples
when the likelihood function is costly to evaluate but easy to sample from. Thus, ABC is useful
in applications where it is possible to simulate data from given model parameters, but computing the likelihood is intractable. ABC was first introduced in population genetics, where one is
interested in recovering the history of a population from genetic data. The history is represented
by a genealogy tree, which is easy to sample from given parameters such as mutation and growth
rates, but the likelihood of the corresponding model is intractable. Likelihood-free methods that
avoid evaluating the likelihood are important in many applications beyond population genetics,
including epidemiology and cosmology.

The key idea behind ABC is to replace the evaluation of likelihoods with a comparison between sampled and observed data. This comparison is based on a choice of distance and tolerance. In its original, most basic form, ABC is the following accept/reject method.

\begin{algorithm}[H]
\caption{ABC Rejection Sampler}
\label{alg:2.4}
\begin{algorithmic}[1]
\Require Prior $f(\theta)$, likelihood $f(y|\theta)$, distance $d$, tolerance $\varepsilon$, data $y_o$, sample size $N$.
\State \textbf{For} $n = 1, \ldots, N$ \textbf{do:}
\State \quad \textbf{Step 1:} Sample $\theta^{(n)} \sim f(\theta)$ from the prior.
\State \quad \textbf{Step 2:} Sample data $y^{(n)} \sim f(y|\theta^{(n)})$ from the likelihood.
\State \quad If $d(y_o, y^{(n)}) \le \varepsilon$, accept $\theta^{(n)}$. Otherwise, go to step~1.
\Ensure $\{\theta^{(n)}\}_{n=1}^N \iid f(\theta|d(y_o, y_{\mathrm{sim}}) \le \varepsilon)$, where $y_{\mathrm{sim}} \sim f(y)$.
\end{algorithmic}
\end{algorithm}

\medskip
\noindent\fbox{\parbox{\dimexpr\linewidth-2\fboxsep-2\fboxrule}{%
\textbf{Pros and Cons}: ABC is useful when the likelihood is easy to sample from but expensive to evaluate. For some Bayesian inference problems, ABC is virtually the only implementable method. Some caveats of ABC include that (i) it only produces approximate posterior samples; (ii) theoretical developments are scarce; and (iii) the tolerance and the distance are often chosen heuristically, and the results are sensitive to these choices.}}
\medskip

\paragraph{Choice of Distance}
It is often impractical to define a distance $d(y_o, y^{(n)})$ between the full
data-sets. Instead, it is preferable to use lower-dimensional sufficient summary statistics $S(y_o)$
and $S(y^{(n)})$ of the data-sets and define
\[
  d(y_o, y^{(n)}) = \tilde{d}\bigl(S(y_o), S(y^{(n)})\bigr),
\]
where $\tilde{d}$ is a distance function defined on the summary statistic space.

\paragraph{Choice of Tolerance}
In choosing $\varepsilon$, there is a trade-off between computational burden and
approximation of the posterior. Smaller $\varepsilon$ decreases the acceptance probability, but leads to
samples whose distribution is closer to the posterior. This trade-off will be illustrated in the
following example.

\begin{example}[ABC: Beta-Binomial Bayesian Model]
\label{ex:2.8}
Consider a random experiment which
involves $n_o$ flips of a coin with unknown probability of heads $p$. Letting $y$ denote the number of
heads observed out of the $n_o$ flips, we have a $\Binom(n_o, p)$ likelihood function
\[
  f(y|p) = \binom{n_o}{y} p^y (1-p)^{n_o - y}.
\]
Taking a $\Beta(\alpha, \beta)$ prior for $p$, with p.d.f.\ $f(p) \propto p^{\alpha-1}(1-p)^{\beta-1}$, $p \in (0,1)$, leads to a
posterior distribution
\begin{align*}
  f(p|y) &\propto f(y|p) f(p) \\
         &\propto p^{\alpha+y-1}(1-p)^{n_o+\beta-y-1}, \quad 0 < p < 1.
\end{align*}
That is, the posterior of $p$ given $y$ is a $\Beta(\alpha + y, n_o + \beta - y)$ distribution. For the purpose of
this example, we choose $\alpha = \beta = 1$, which corresponds to a Unif$(0,1)$ prior.

Next, for the sake of the example, suppose that we do not know the true posterior distribution
of $p$. We will approximate it using ABC. To that end, we simulate the model using a sample
size of $n_o = 500$ with $p = 0.4$ and record the number $y_o$ of observed successes. We seek
to approximate the posterior of $p$ given $y_o$. We repeatedly sample values $p^{(n)}$ from the prior
distribution $\Beta(1,1)$, and simulate the number of successes $y^{(n)} \sim f(\cdot|p^{(n)})$. We next define a
distance function,
\[
  d(y_o, y^{(n)}) := \frac{1}{n_o} \cdot \abs{y^{(n)} - y_o}.
\]
Using three tolerance values $\varepsilon = 0.01, 0.02, 0.05$, we accept the sample $p^{(n)}$ if $d(y^{(n)}, y_o) < \varepsilon$.
For each choice of tolerance, we continue this procedure until we obtain $N = 1000$ accepted
samples. Finally, for each tolerance, we plot a histogram of the accepted samples against the true
posterior distribution, which is $\Beta(1 + y_o, n_o + 1 - y_o)$. Figure~\ref{fig:2.4} shows the results, along
with the number of samples generated in each case to reach $N = 1000$ accepted samples.
\end{example}

\begin{figure}[htbp]
\centering
\includegraphics[width=0.75\textwidth]{figures/fig_02_4}
\caption{\textit{ABC for a Beta-binomial model as described in Example~\ref{ex:2.8}. The results illustrate
the trade-off between cost and accuracy in the choice of tolerance.}}
\label{fig:2.4}
\end{figure}
